% Generated draft: requires mathematical and source review.
\documentclass[11pt]{article}
\usepackage[margin=25mm]{geometry}
\usepackage{fontspec}
\setmainfont{Latin Modern Roman}
\usepackage{amsmath,amssymb,mathtools,unicode-math}
\setmathfont{Latin Modern Math}
\usepackage{enumitem,needspace,xurl,hyperref}
\setlength{\parindent}{0pt}
\setlength{\parskip}{5pt}
\setlength{\emergencystretch}{3em}
\newcommand{\N}{\mathbb{N}}
\newcommand{\Z}{\mathbb{Z}}
\newcommand{\Q}{\mathbb{Q}}
\newcommand{\R}{\mathbb{R}}
\newcommand{\C}{\mathbb{C}}
\newcommand{\F}{\mathbb{F}}
\newcommand{\A}{\mathbb{A}}
\newcommand{\PP}{\mathbb{P}}
\newcommand{\E}{\mathbb{E}}
\newcommand{\eps}{\varepsilon}
\newcommand{\dd}{\,\mathrm{d}}
\begin{document}
\section*{UnsolvedMath ID: 30004289}
\textbf{Source problem code:} catalogue entry
\subsection{Definitions and mathematical statement}
For a fundamental discriminant \(d\), let \(\chi_d\) denote its associated quadratic Dirichlet character. The supplied record concerns the finite Dirichlet polynomial
\[
P_A(X;d)=\sum_{1\leq n\leq X}\frac{\tau_A(n)\chi_d(n)}{\sqrt n}.
\]
Here \(A\) and the coefficients \(\tau_A(n)\) are notation retained from the source; the record does not define them or specify their admissible values.

The proposed long-polynomial recipe asserts, schematically, that
\[
\text{a smoothed average of }P_A(X;d)\text{ over fundamental discriminants }d
=
\text{an integral of quadratic }L\text{-function swap terms}
+O\!\left(D^{1-\delta}\right)
\qquad(D\to\infty).
\]
This display records the advertised target, not a fully specified mathematical conjecture. The supplied record omits the averaging weight and normalization, the discriminant range and sign restrictions, the relation between \(X\) and \(D\), the definition and admissible shifts of \(A\), the swap terms, and the integral, including its contour or domain. It also omits the quantification and parameter dependence of \(\delta\), the uniformity of the error estimate, and the permitted dependence of its implied constant. A positive \(\delta\) would express a power saving, but the required hypotheses for such a saving are not supplied. These omissions prevent an exact quantified formulation.

The supplied dated literature assessment describes a conditional result of Conrey--Rodgers in a polynomial-length range with exponent below \(2\), corresponding to the \(0\)- and \(1\)-swap terms. It describes the longer range and the full shifted-moment prediction as remaining open. The conditioning hypotheses, definition of that exponent, and full shifted-moment formula are absent from the record; this assessment is retained as supplied curation, not as an independently verified theorem or status determination.
\subsection{Short English statement}
The proposed recipe predicts a smoothed average of quadratic-character Dirichlet polynomials from an integral of swap terms, with error \(O(D^{1-\delta})\). The supplied record omits the formulas and hypotheses needed to make this precise. Its dated curation labels the problem partially solved and reports a conditional result for a length exponent below \(2\); these claims require independent mathematical review.
\subsection{Sources}
\begin{itemize}
\item Supplied local UnsolvedMath registry record: permanent ID \(30004289\), title ``Long Dirichlet-Polynomial Recipe for Quadratic \(L\)-Functions,'' exact source problem code \texttt{catalogue entry}, in the Oberwolfach Reports Open Problems collection. The record gives proposed year \(2019\) and names no proposer.
\item Supplied original report citation: \emph{Analytic Number Theory (2020)}, Oberwolfach Reports, pp.~3141--3205, DOI \texttt{10.4171/owr/2019/50}. \url{https://doi.org/10.4171/owr/2019/50}. The supplied background describes these reports as documenting mathematical workshops and their open-problem discussions. The relevant problem passage and full recipe are not reproduced in the record.
\item Supplied literature reference: Conrey--Rodgers, ``quadratic twists of long Dirichlet polynomials.'' \url{https://doi.org/10.1007/s00209-025-03848-x}.
\item Supplied literature reference: Conrey--Rodgers, ``conditional theorem.'' \url{https://arxiv.org/abs/2208.01783}.
\item Literature assessment date: \texttt{2026-08-22}. This is supplied dated curation, not independently verified. It labels the problem partially solved, reports a conditional range with polynomial-length exponent below \(2\) corresponding to the \(0\)- and \(1\)-swap terms, and describes the longer range and full shifted-moment prediction as remaining open. The supplied background explicitly requests independent verification of time-sensitive status and cited claims.
\end{itemize}
\end{document}
